\section{Conclusion}

As AI inference scales to models that far exceed accelerator memory, the
efficiency of hierarchical memory orchestration increasingly determines the
economic and environmental cost of deployed AI.
This work demonstrates that such orchestration is not a continuously
optimisable process: it exhibits intrinsic, regime-dependent behaviour governed
by two dimensionless ratios---fast-memory residency and cross-tier
transfer-pressure---whose critical thresholds define abrupt, phase-like
transitions between operationally distinct regimes.

By linking these regimes to dominant terms in a minimal multi-tier transfer
model, we derive analytically interpretable boundaries that explain \emph{why}
orchestration strategies effective in coordination-dominated regimes fail
predictably elsewhere---and show that such failures reflect physical constraints
rather than implementation shortcomings.
Extensive validation across language and vision workloads on five hardware
platforms confirms that regime boundaries are hardware-agnostic and recoverable
to within $\pm 0.05$ of predicted thresholds.

The regime framework reframes hierarchical memory optimisation from a problem
of continuous tuning to one of regime-aware design: the first question for any
memory-bound inference system is not ``how much more coordination?'' but ``in
which regime does the system operate?''
This perspective applies beyond large language models to any computational
workload whose active working set cannot fully reside in fast memory---including
multimodal inference, retrieval-augmented generation, data-intensive analytics,
and HPC-style scientific computing---and provides a principled foundation for
understanding scalability, stability, and energy efficiency in increasingly
complex memory hierarchies.
